\documentclass[11pt,
%10pt, 
%hyperref={colorlinks},
aspectratio=169,
hyperref={colorlinks}
]{beamer}
\usetheme{Singapore}
\usecolortheme[snowy, cautious]{owl}

\usepackage[utf8]{inputenc}
\usepackage[T1]{fontenc}
\usepackage[american]{babel}
\usepackage{graphicx}
\usepackage{hyperref}
\hypersetup{
        colorlinks=true,
        urlcolor=[rgb]{1,0,1},
        linkcolor=[rgb]{1,0,1}}

\usepackage[natbib=true,style=authoryear,backend=bibtex,useprefix=true]{biblatex}
\usepackage{listings}
\lstset{numbers=right, 
numberstyle=\tiny, 
%        breaklines=true,
%        backgroundcolor=\color{light-gray},
numbersep=5pt,
xleftmargin=\parindent,
xrightmargin=.25in} 

%\setbeamercolor*{bibliography entry title}{fg=black}
%\setbeamercolor*{bibliography entry location}{fg=black}
%\setbeamercolor*{bibliography entry note}{fg=black}
\definecolor{OwlGreen}{RGB}{75,0,130} % easier to see
\setbeamertemplate{bibliography item}{}
\setbeamerfont{caption}{size=\footnotesize}
\setbeamertemplate{frametitle continuation}{}
\setcounter{tocdepth}{1}
\renewcommand*{\bibfont}{\scriptsize}
\addbibresource{bibliography.bib}

\renewcommand*{\thefootnote}{\fnsymbol{footnote}}

%\author{\copyright\hspace{1pt}Ashrith Barthur\footnote{\tiny{This material is shared under a \href{https://creativecommons.org/licenses/by/4.0/deed.ast}{CC By 4.0 license} which allows for editing and redistribution, even for commercial purposes. However, any derivative work should attribute the author and H2O.AI.}}}
\author{Ashrith Barthur}
\title{How to Write A Recipe, and a TimeSeries Recipe?}
\subtitle{Automating Feature Engineering Using DriverlessAI}
\logo{\includegraphics[height=8pt]{img/h2o_logo.png}}
\institute{\href{https://www.h2o.ai}{H\textsubscript{2}O.ai}}
\date{\today}
\subject{How to Write A Recipe?}

\begin{document}

\maketitle

%	\begin{frame}

%		\frametitle{Contents}

%		\tableofcontents{}

%	\end{frame}

%-------------------------------------------------------------------------------
\subsection{Question}
%-------------------------------------------------------------------------------
\begin{frame}
        \frametitle{Question}
        \begin{enumerate}
                \item How many of us have built variables, features, transformers, or feature transformers?
                \item What are they?
        \end{enumerate}
\end{frame}
%-------------------------------------------------------------------------------
\subsection{Answer}
%-------------------------------------------------------------------------------
\begin{frame}
        \frametitle{Answer}
        \begin{enumerate}
                \item Variables, features, transformers, feature transformers all refer to the same.
                \item Each column in your data is considered a variable(\textit{incoming}) or a feature(\textit{incoming}).
                \item Each \textit{new} column created is also referred to as a variable or a feature.
                \item The process of creating a new variable, or a feature is called a transformation. 
                \item The code processing an \textit{existing} column to a \textit{new} column is called a \textit{transformer}.
        \end{enumerate}
\end{frame}
%-------------------------------------------------------------------------------
\subsection{Example Transformation}
%-------------------------------------------------------------------------------
\begin{frame}[fragile]
        \frametitle{Example Transformation}
        \begin{enumerate}
                \item \textit{height}- Variable
                \item New variable after transformation \textit{log2(height)}
        \end{enumerate}
\end{frame}
%-------------------------------------------------------------------------------
\subsection{Question}
%-------------------------------------------------------------------------------
\begin{frame}
        \frametitle{Question}
        \begin{enumerate}
                \item How many of us are familiar with Custom Transformers in Driverless AI?
                \item What are they?
        \end{enumerate}
\end{frame}
%-------------------------------------------------------------------------------
\subsection{Answer}
%-------------------------------------------------------------------------------
\begin{frame}
        \frametitle{Answer}
        \begin{enumerate}
                \item DriverlessAI already has a large, comprehensive set of transformers. 
                \item But there are always domains that require nuanced features. 
                \item And for this, DriverlessAI provides us to create custom transformers. 
                \item This is provided by provisioning an extension class \textit{CustomTransformer}
        \end{enumerate}
\end{frame}
%-------------------------------------------------------------------------------
\subsection{Recipe Introduction}
%-------------------------------------------------------------------------------
\begin{frame}[fragile]
        \frametitle{How Did We Build A Custom Transformer?}
        Driverless AI provides an extension. \\
        This is a class `CustomTransformer`
        \begin{verbatim}
            class ExampleLogTransformer(CustomTransformer):
        \end{verbatim}
\end{frame}
%-------------------------------------------------------------------------------
\section{Feature Recipe Structure}
%-------------------------------------------------------------------------------
\begin{frame}[fragile]
        \frametitle{How Did We Build This?}
        The class has:
        \begin{enumerate}
                \item Parameters that need to be provided. 
                \item These parameters are specific to the type of feature recipe that you are building. 
                \item It also has four methods which primarily handle your feature engineering transformation. 
        \end{enumerate}
\end{frame}
%-------------------------------------------------------------------------------
\subsection{Parameters - Basic}
%-------------------------------------------------------------------------------
\begin{frame}[fragile]
        \frametitle{Parameters - Basic}
        \begin{verbatim}
        class ExampleLogTransformer(CustomTransformer):
            _regression = True
            _binary = True
            _multiclass = True
        \end{verbatim}
\end{frame}
%-------------------------------------------------------------------------------
\subsection{Parameters - Advanced}
%-------------------------------------------------------------------------------
\begin{frame}[fragile]
        \frametitle{Parameters - Advanced}
        \begin{verbatim}
        class ExampleLogTransformer(CustomTransformer):
            _regression = True
            _binary = True
            _multiclass = True
            _numeric_output = True
            _is_reproducible = True
            _excluded_model_classes = [`tensorflow']
            _modules_needed_by_name = ["custom_package==1.0.0"]
        \end{verbatim}
\end{frame}
%-------------------------------------------------------------------------------
%		\subsection{Parameters - Advanced}
%-------------------------------------------------------------------------------
\begin{frame}[fragile]
        \frametitle{Acceptance Method}
        \begin{verbatim}
        class ExampleLogTransformer(CustomTransformer):
            _regression = True
            _binary = True
            _multiclass = True
            _numeric_output = True
            _is_reproducible = True
            _excluded_model_classes = [`tensorflow']
            _modules_needed_by_name = ["custom_package==1.0.0"]

            @staticmethod
            def do_acceptance_test():
                return True
        \end{verbatim}
\end{frame}
%-------------------------------------------------------------------------------
%		\subsection{Parameters - Advanced}
%-------------------------------------------------------------------------------
\begin{frame}[fragile]
        \frametitle{Input Data}
        \begin{verbatim}
            ...
            @staticmethod
            def do_acceptance_test():
                return True

            @staticmethod
            def get_default_properties():
                return dict(col_type = "numeric"
                    ,min_cols = 1, max_cols = 1, 
                    relative_importance = 1)
        \end{verbatim}
\end{frame}

%-------------------------------------------------------------------------------
%		\subsection{Parameters - Advanced}
%-------------------------------------------------------------------------------
\begin{frame}[fragile]
        \frametitle{Input Data Types}
        \begin{verbatim}
        a. "all"    - all column types
        b. "any"    - any column types
        c. "numeric"    - numeric int/float column
        d. "categorical" - string/int/float column considered a categorical 
                            for feature engineering
        e. "numcat" - allow both numeric or categorical
        f. "datetime"    - string or int column with raw datetime such as 
                            `%Y/%m/%d %H:%M:%S' or `%Y%m%d%H%M'
        \end{verbatim}
\end{frame}
%-------------------------------------------------------------------------------
%		\subsection{Parameters - Advanced}
%-------------------------------------------------------------------------------
\begin{frame}[fragile]
        \frametitle{Input Data Types}
        \begin{verbatim}
        g. "date"   - string or int column with raw date such as 
                      `%Y/%m/%d' or `%Y%m%d'
        h. "text"   - string column containing text 
                       (and hence not treated as categorical)
        i. "time_column" - the time column specified at the start of 
                            the experiment (unmodified)
        \end{verbatim}
\end{frame}

%-------------------------------------------------------------------------------
%		\subsection{Parameters - Advanced}
%-------------------------------------------------------------------------------
\begin{frame}[fragile]
        \frametitle{Fit Function}
        \begin{verbatim}
            @staticmethod
            def get_default_properties():
                return dict(col_type = "numeric"
                    ,min_cols = 1, max_cols = 1,
                    relative_importance = 1)

            def fit_transform(self, X: dt.Frame, y: np.array = None):
                X_pandas = X.to_pandas()
                X_p_log = np.log10(X_pandas)
                return X_p_log
        \end{verbatim}
\end{frame}
%-------------------------------------------------------------------------------
%		\subsection{Parameters - Advanced}
%-------------------------------------------------------------------------------
\begin{frame}[fragile]
        \frametitle{Transform Function}
        \begin{verbatim}
        def fit_transform(self, X: dt.Frame, y: np.array = None):
            X_pandas = X.to_pandas()
            X_p_log = np.log10(X_pandas)
            return X_p_log

        def transform(self, X: dt.Frame):
            X_pandas = X.to_pandas()
            X_p_log = np.log10(X_pandas)
            return X_p_log
        \end{verbatim}
\end{frame}
%-------------------------------------------------------------------------------
%		\subsection{Parameters - Advanced}
%-------------------------------------------------------------------------------
\begin{frame}[fragile]
        \frametitle{Library}
        \begin{verbatim}
        from h2oaicore.systemutils import segfault, 
            ,loggerinfo, main_logger
        from h2oaicore.transformer_utils 
            import CustomTransformer
        import datatable as dt
        import numpy as np
        import pandas as pd
        import logging
        \end{verbatim}
\end{frame}
%-------------------------------------------------------------------------------
%		\subsection{Parameters - Advanced}
%-------------------------------------------------------------------------------
\begin{frame}[fragile]
        \frametitle{Time Series Introduction Auto Arima}
        \begin{enumerate}
                \item In our example we will bring in the \textit{auto\_arima}  function as a part of the recipe. 
                \item This is available in the \textit{pmdarima} package available for \textit{Python}.
                \item The \textit{auto\_arima} function tries different `p`, `q`, and `d` values for \textit{ARIMA}, automatically. 
                \item It selects the best values based on the lowest value in the information criterion. 
        \end{enumerate}
\end{frame}
%-------------------------------------------------------------------------------
%		\subsection{Parameters - Advanced}
%-------------------------------------------------------------------------------
\begin{frame}[fragile]
        \frametitle{Recap on ARIMA}
        What is ARIMA?\\
        ARIMA stands for Auto Regressive Integrated Moving Average
\end{frame}
%-------------------------------------------------------------------------------
%		\subsection{Parameters - Advanced}
%-------------------------------------------------------------------------------
\begin{frame}[fragile]
        \frametitle{Recap on ARIMA}
        Auto Regressive means the target depends on its own lags:
        \begin{enumerate}
                \item $Y_{t} = \alpha  + \beta_{1}Y_{t-1}+...+\beta_{p}Y_{t-p}+\epsilon_{t}$
        \end{enumerate}
        Integrated means that the target has been differenced to make the time series stationary:
        \begin{enumerate}
                \item First order differencing: $Y_{t}^{d1} = Y_{t} - Y_{t-1}$
                \item Second order differencing: $Y_{t}^{d2} = Y_{t}^{d1} - Y_{t-1}^{d1}$
        \end{enumerate}
        Moving Average means the target depends on previous prediction errors:
        \begin{enumerate}
                \item $Y_{t} = \alpha + \phi_{1}\epsilon_{t-1} + ... + \beta_{q}\epsilon_{t-q}$
        \end{enumerate}

\end{frame}
%-------------------------------------------------------------------------------
\section{Feature Recipe Structure}
%-------------------------------------------------------------------------------
\begin{frame}[fragile]
        \frametitle{TimeSeries Recipe Basics}
        \begin{enumerate}
                \item There is a custom class for creating TimeSeries recipes \textit{CustomTimeSeriesTransformer}. 
                \item Similar to \textit{CustomTransformer}, \textit{CustomTimeSeriesTransformer} has pre-defined parameters and functions. 
        \end{enumerate}

\end{frame}
%-------------------------------------------------------------------------------
\subsection{Parameters - Basic}
%-------------------------------------------------------------------------------
\begin{frame}[fragile]
        \frametitle{Generic Recipe Parameters}
        \begin{verbatim}
        class MyAutoArimaTransformer(CustomTimeSeriesTransformer):
            _binary = False
            _multiclass = False
            _modules_needed_by_name = [`pmdarima']
            _included_model_classes = None
        \end{verbatim}

\end{frame}
%-------------------------------------------------------------------------------
%		\subsection{Question}
%-------------------------------------------------------------------------------
\begin{frame}[fragile]
        \frametitle{TimeSeries Recipe Specific Parameters}
        \begin{verbatim}
            self.tgc = kwargs[`tgc']  
            self.target = kwargs[`target']  
            if isinstance(kwargs[`time_column'],list):
                self.time_column = kwargs[`time_column'][0] 
            else:
                self.time_column kwargs[`time_column']
        \end{verbatim}
\end{frame}
%-------------------------------------------------------------------------------
\section{Feature Recipe Structure}
%-------------------------------------------------------------------------------
\begin{frame}[fragile]
        \frametitle{TimeSeries Recipe Specific Parameters}
        There are three parameters primarily required by CustomTimeSeries class.
        \begin{enumerate}
                \item \verb|self.tgc| Time series groups
                \item \verb|self.target| The target column
                \item \verb|self.time_column| The column that holds time. 
        \end{enumerate}

\end{frame}
%-------------------------------------------------------------------------------
\section{Feature Recipe Structure}
%-------------------------------------------------------------------------------
\begin{frame}[fragile]
        \frametitle{TimeSeries Recipe Class}
        The \textit{CustomTimeSeriesTransformer} class shares the basic, four methods of \textit{CustomTransformer} Class. These are methods that DriverlessAI invokes while running custom recipes.
        \begin{enumerate}
                \item \verb|do_acceptance_test|
                \item \verb|get_default_properties|
                \item \verb|fit_transform|
                \item \verb|transform|
        \end{enumerate}

\end{frame}
%-------------------------------------------------------------------------------
\section{Feature Recipe Structure}
%-------------------------------------------------------------------------------
\begin{frame}[fragile]
        \frametitle{TimeSeries Recipe Class}
        Additionally, there are two other functions that are invokable in \textit{CustomTimeSeriesTransformer} class. They are:
        \begin{enumerate}
                \item \verb|fit| builds the model to which the data will fit. 
                \item \verb|update_history| updates the model fit with additional data. 
        \end{enumerate}

\end{frame}
%-------------------------------------------------------------------------------
\subsection{Parameters - Basic}
%-------------------------------------------------------------------------------
\begin{frame}[fragile]
        \frametitle{Acceptance and Properties Methods}
        \begin{verbatim}
        @staticmethod
        def do_acceptance_test():
        return False

        @staticmethod
        def get_default_properties():
            return dict(col_type="time_column"
                ,min_cols=1, max_cols=1,
                relative_importance=1)
        \end{verbatim}
\end{frame}
%-------------------------------------------------------------------------------
%		\subsection{Question}
%-------------------------------------------------------------------------------
\begin{frame}[fragile]
        \frametitle{Building the Fit Function}
        \begin{verbatim}
        def fit(self, X: dt.Frame, y: np.array = None):
            pm = importlib.import_module(`pmdarima')
            self.models = {}
            X = X.to_pandas()
            XX = X[self.tgc].copy()
            XX[`y'] = np.array(y)
            self.nan_value = np.mean(y)
            elf.ntrain = X.shape[0]
        \end{verbatim}
\end{frame}
%-------------------------------------------------------------------------------
%		\subsection{Question}
%-------------------------------------------------------------------------------
\begin{frame}[fragile]
        \frametitle{Building the Fit Function}
        \begin{verbatim}
            tgc_wo_time = list(np.setdiff1d(self.tgc, self.time_column))
            if len(tgc_wo_time) > 0:
                XX_grp = XX.groupby(tgc_wo_time)
            else:
                XX_grp = [([None], XX)]
        \end{verbatim}
\end{frame}
%-------------------------------------------------------------------------------
%		\subsection{Question}
%-------------------------------------------------------------------------------
\begin{frame}[fragile]
        \frametitle{Building the Fit Function}
        \begin{verbatim}
            nb_groups = len(XX_grp)
            for _i_g, (key, X) in enumerate(XX_grp):
                key = key if isinstance(key, list) else [key]
                grp_hash = `_'.join(map(str, key))
                order = np.argsort(X[self.time_column])
                try:
                    model = pm.auto_arima(X[`y'].values[order] 
                            ,error_action=`ignore')
                    except:
                    model = None
                self.models[grp_hash] = model
                return self
        \end{verbatim}
\end{frame}
%-------------------------------------------------------------------------------
%		\subsection{Question}
%-------------------------------------------------------------------------------
\begin{frame}[fragile]
        \frametitle{Building the Transform Function}
        \begin{verbatim}
            nb_groups = len(XX_grp)
            preds = []
            for _i_g, (key, X) in enumerate(XX_grp):
                key = key if isinstance(key, list) else [key]
                grp_hash = `_'.join(map(str, key))
                order = np.argsort(X[self.time_column])
        \end{verbatim}
\end{frame}
%-------------------------------------------------------------------------------
%		\subsection{Question}
%-------------------------------------------------------------------------------
\begin{frame}[fragile]
        \frametitle{Building the Transform Function}
        \begin{verbatim}
            if grp_hash in self.models:
	            model = self.models[grp_hash]
	            if model is not None:
	                if hasattr(self, `is_train'): 
	                    yhat = model.predict_in_sample()
	                else:
	                    model.predict(n_periods=X.shape[0])
	                yhat = yhat[order]
	                XX = pd.DataFrame(yhat, columns=[`yhat'])
	             else:
	                XX = pd.DataFrame(np.full((X.shape[0], 1), self.nan_value) 
	                     ,columns=[`yhat'])  # invalid model
	                        ...
        \end{verbatim}
\end{frame}
%-------------------------------------------------------------------------------
%		\subsection{Question}
%-------------------------------------------------------------------------------
\begin{frame}[fragile]
        \frametitle{Building the Transform Function}
        \begin{verbatim}
                            ...
            else:
                XX = pd.DataFrame(np.full((X.shape[0], 1), self.nan_value), 
                columns=[`yhat'])  # unseen groups
                XX.index = X.index
                preds.append(XX)
            XX = pd.concat(tuple(preds), axis=0).sort_index()
            return XX
        \end{verbatim}
\end{frame}
%-------------------------------------------------------------------------------
%		\subsection{Question}
%-------------------------------------------------------------------------------
\begin{frame}[fragile]
        \frametitle{Building the Fit Transform Function}
        \begin{verbatim}
            def fit_transform(self, X: dt.Frame, y: np.array = None):
                self.is_train = True
                ret = self.fit(X, y).transform(X)
                del self.is_train
                return ret
        \end{verbatim}
\end{frame}
%-------------------------------------------------------------------------------
%		\subsection{Question}
%-------------------------------------------------------------------------------
\begin{frame}[fragile]
        \frametitle{Building the Update History Function}
        \begin{verbatim}
            def update_history(self, X: dt.Frame, y: np.array = None):
                X = X.to_pandas()
                XX = X[self.tgc].copy()
                XX[`y'] = np.array(y)
                tgc_wo_time = list(np.setdiff1d(self.tgc, self.time_column))
                if len(tgc_wo_time) > 0:
                   XX_grp = XX.groupby(tgc_wo_time)
                else:
                   XX_grp = [([None], XX)]
        \end{verbatim}
\end{frame}
%-------------------------------------------------------------------------------
%		\subsection{Question}
%-------------------------------------------------------------------------------
\begin{frame}[fragile]
        \frametitle{Building the Update History Function}
        \begin{verbatim}
                for key, X in XX_grp:
                    key = key if isinstance(key, list) else [key]
                    grp_hash = `_'.join(map(str, key))
                    order = np.argsort(X[self.time_column])
                    if grp_hash in self.models:
                        model = self.models[grp_hash]
                         if model is not None:
                            model.update(X[`y'].values[order])
                return self
        \end{verbatim}
\end{frame}
%-------------------------------------------------------------------------------
%		\subsection{Parameters - Advanced}
%-------------------------------------------------------------------------------
\begin{frame}[fragile]
        \frametitle{Library}
        \begin{verbatim}
        from h2oaicore.systemutils import segfault, 
            ,loggerinfo, main_logger
        from h2oaicore.transformer_utils 
            import CustomTimeSeriesTransformer
        import datatable as dt
        import numpy as np
        import pandas as pd
        import logging
        \end{verbatim}
\end{frame}
%-------------------------------------------------------------------------------
%\subsection{Question}
%-------------------------------------------------------------------------------
\begin{frame}
        \frametitle{Advantages}
        \begin{enumerate}
                \item Feature engineering process standardised by:
                        \begin{enumerate}
                                \item preset parameters
                                \item preset methods
                        \end{enumerate}
                \item Effort minimisation leads to minimisation in time spent.
                \item Build only once - Feature engineering is carried over from training/testing to production.
                \item DAI automatically, runs multiple models on various sets of features to get the best model. 
                \item All the requirements are handled internally by DAI. 
        \end{enumerate}
\end{frame}
%-------------------------------------------------------------------------------
%	References
%-------------------------------------------------------------------------------

\begin{frame}[t, allowframebreaks]

        \frametitle{References}	

        \textbf{How to build a recipe}\\
        \small{\url{https://github.com/ashrith/how_to_write_a_recipe}}

        \framebreak		

        \printbibliography

\end{frame}
%-------------------------------------------------------------------------------
\section{Questions}
%------------------------------------------------------------------------------

\begin{frame}

        \frametitle{Thanks \& Questions}
        \begin{itemize}
                \item Olivier Grellier, Ph.D, Data Scientist, Kaggle Grandmaster
        \end{itemize}

\end{frame}

%-------------------------------------------------------------------------------
%	\section{End}
%------------------------------------------------------------------------------
\end{document}
